% Propietario conservado: /Users/ruben/Documents/New project/output/REVISION_ARGUMENTAL_APERTURAS_CIERRES_20260915/VI/delivery/MANUSCRIPT_VI_ES_EN_COMPLETE/source_es/sections/06_kepler_sommerfeld.tex
\section{Central realisation of trajectories and Sommerfeld holonomy}
\label{vi:vi:sec:kepler}

The particle defined in section~\ref{vi:vi:sec:particula} has
an observable history before receiving an orbital realization.
The question of this section is how a representation of that
history preserves flux, area and phase when expressed as
central motion. The answer is organized into four operations:
cancellation of the interior fluxes in a graph cut,
normalization of its boundary in a metric chart, evolution under
the declared central operator and phase transport over its cycles.

The order keeps visible which datum enters each statement.
The graph and its records come from APP--TRIT--TPK; the metric
chart specifies how its boundary is measured; the dynamical realization
fixes the Lagrangian acting in that chart; the semiclassical condition
uses the reduced Planck constant already constructed and the holonomy
of the transported line. The results of orbital geometry are
proved with these explicit starting data. Central
mechanics is recognized in the resulting realization, without making its
conventional constants inputs to the generator.

\subsection{Flux conservation and boundary reading}
\label{vi:vi:subsec:gauss}

Let $G=(V,E)$ be a finite cut of the route graph of the
realisation under consideration, with an arc in each orientation of its
edges. An oriented flux is a function
$F:\vec E\to\mathbb R$ satisfying
$F(y,x)=-F(x,y)$. At a vertex $x$, its divergence is
\begin{equation}
 \operatorname{div}F(x)
 =\sum_{y:\{x,y\}\in E}F(x,y).
 \label{vi:vi:eq:divergencia-grafo}
\end{equation}
For $S\subseteq V$, the outward flux is evaluated over the
arcs originating in $S$ and ending in $V\setminus S$.

\begin{theorem}[Flux balance in a finite cut]
\label{vi:vi:thm:gauss}
For every subset $S\subseteq V$,
\begin{equation}
 \sum_{x\in S}\operatorname{div}F(x)
 =\sum_{\substack{x\in S,\ y\notin S\\\{x,y\}\in E}}F(x,y).
 \label{vi:vi:eq:gauss}
\end{equation}
\end{theorem}
\begin{proof}
An edge with both endpoints in $S$ contributes twice to the
sum of divergences, through $F(x,y)+F(y,x)=0$. Edges
with both endpoints outside $S$ do not contribute. Exactly
the edges crossing the boundary remain, with the
indicated orientation.
\end{proof}

The theorem transmits to the boundary the balance produced in the
interior. For a cubic representation with nearest neighbours,
consider
\[
 C_R=\{x\in\mathbb Z^3:\|x\|_\infty\leq R\},
 \qquad R\in\mathbb Z_{\geq0}.
\]
Each of the six faces contains $(2R+1)^2$ outward
links in its normal direction. Links in different
directions are counted separately, even when they originate at
the same edge or corner vertex. Hence,
\begin{equation}
 |\partial C_R|=6(2R+1)^2.
 \label{vi:vi:eq:frontera-cubica}
\end{equation}
This census counts outward links, not the number of
boundary vertices.

If $Q_F$ is the total flux of the enclosed source, its average
per link is $Q_F/[6(2R+1)^2]$. Under the condition of mean isotropy
of the reader, that average is the intensity assigned to the
shell. Let us define its effective radius, in units of the step of
the chart, by
\begin{equation}
 4\pi r_{\mathrm{eff}}^2=6(2R+1)^2.
 \label{vi:vi:eq:radio-efectivo}
\end{equation}
We then obtain
\begin{equation}
 E_R=\frac{Q_F}{4\pi r_{\mathrm{eff}}^2}.
 \label{vi:vi:eq:inverso-area}
\end{equation}
The coordinate $\pi$ of this normalization is the regional
HMT output already published. The identity follows from the census and the
definition of the radius; the passage from an average to a central
realization additionally uses the stated isotropy. The dimensional
chart then decides whether $Q_F$ represents electric,
gravitational or other observable flux and how it converts the lattice
step into length. The finite balance preserves this
distinct status of each operation.

\subsection{Central dynamics on the Euclidean sheet}
\label{vi:vi:subsec:dinamica-central}

Let $\mu>0$ be the reduced mass of the material realisation and let
$\mathcal K>0$ be its central coupling, with the dimension of
energy times length. These coordinates are received after
the internal constructions of mass, action and coupling;
their origin and dimensional line are preserved when carrying out the
realisation. The domain of the following theorem consists of all
realisations that adopt, on their Euclidean sheet,
\begin{equation}
 L(\mathbf r,\dot{\mathbf r})
 =\frac\mu2\|\dot{\mathbf r}\|^2+\frac{\mathcal K}{r},
 \qquad
 H(\mathbf r,\mathbf p)
 =\frac{\|\mathbf p\|^2}{2\mu}-\frac{\mathcal K}{r},
 \quad r=\|\mathbf r\|>0,
 \quad\mathbf p=\mu\dot{\mathbf r}.
 \label{vi:vi:eq:kepler-lh}
\end{equation}
The exclusion of $r=0$ delimits the regular domain of these
coordinates. The trajectory is considered over a collision-free
time interval.

\begin{proposition}[Orbital plane and conic equation]
\label{vi:vi:prop:conica}
A non-radial solution of the Euler--Lagrange equation of
\eqref{vi:vi:eq:kepler-lh} conserves
$\boldsymbol\ell=\mathbf r\times\dot{\mathbf r}$, with
$\ell=\|\boldsymbol\ell\|>0$. In the plane orthogonal to
$\boldsymbol\ell$, it can be written
\begin{equation}
 r(\theta)=\frac{p}{1+\epsilon\cos(\theta-\theta_0)},
 \qquad p=\frac{\mu\ell^2}{\mathcal K},
 \qquad \epsilon\geq0.
 \label{vi:vi:eq:conica}
\end{equation}
Negative energy corresponds to the elliptic regime
$0\leq\epsilon<1$.
\end{proposition}
\begin{proof}
Differentiating the potential gives
\begin{equation}
 \mu\ddot{\mathbf r}=-\frac{\mathcal K}{r^3}\mathbf r.
 \label{vi:vi:eq:fuerza-central}
\end{equation}
Hence,
$\dot{\boldsymbol\ell}=\mathbf r\times\ddot{\mathbf r}=0$.
Position and velocity remain orthogonal to the
non-zero constant vector $\boldsymbol\ell$, fixing the plane.
In its polar coordinates, $r^2\dot\theta=\ell$. Introducing
$u=1/r$ into the radial equation gives Binet's equation
\[
 u''(\theta)+u(\theta)=\frac{\mathcal K}{\mu\ell^2}.
\]
Its general solution is
$u=\mathcal K[1+\epsilon\cos(\theta-\theta_0)]/(\mu\ell^2)$,
and inversion proves \eqref{vi:vi:eq:conica} on the domain
where the denominator is positive. Substitution into
the energy gives
$H=\mathcal K(\epsilon^2-1)/(2p)$, which is negative
exactly when $\epsilon<1$.
\end{proof}

Here the letter $\epsilon$ denotes orbital eccentricity,
not Euler's number. The vector $\boldsymbol\ell$ is
specific angular momentum, not the Planck constant.
These notational choices allow the geometric
observables and the internal constants entering their
realisation to be compared without conflation.

\begin{proposition}[Areal conservation and period law]
\label{vi:vi:prop:ley-areal}
The orbit's areal velocity is constant:
\begin{equation}
 \frac{d\mathcal A_{\mathrm{orb}}}{dt}=\frac\ell2.
 \label{vi:vi:eq:ley-areal}
\end{equation}
For an ellipse with semimajor axis $a$, its period satisfies
\begin{equation}
 T^2=\frac{4\pi^2\mu}{\mathcal K}a^3.
 \label{vi:vi:eq:tercera-ley}
\end{equation}
\end{proposition}
\begin{proof}
The sector swept out in time $dt$ has area
$r^2d\theta/2=\ell\,dt/2$. If $b=a\sqrt{1-\epsilon^2}$
is the semiminor axis, the total area is $\pi ab$ and
$T=2\pi ab/\ell$. The relation
$p=a(1-\epsilon^2)=\mu\ell^2/\mathcal K$ implies
$\ell^2=\mathcal K a(1-\epsilon^2)/\mu$.
Substitution into $T^2=4\pi^2a^2b^2/\ell^2$ gives
\eqref{vi:vi:eq:tercera-ley}.
\end{proof}

Areal conservation is determined by the central symmetry
of the motion. The ellipse of this proposition lives in
configuration space. The action ellipse constructed
from the angular pair, in contrast, lives in its action chart
and transports the corresponding normalizations. The link
between the two requires the realization map that identifies their
variables; sharing an elliptical shape does not replace that map.
The orbital reading thus preserves the dependencies on mass,
coupling and geometry that produce it.

\subsection{Orbital covering and real sign line}
\label{vi:vi:subsec:cubierta-orbital}

On the complex plane with the origin removed, the map
\begin{equation}
 q=z^2:\mathbb C\setminus\{0\}\longrightarrow
                 \mathbb C\setminus\{0\}
 \label{vi:vi:eq:levi-civita}
\end{equation}
is a double covering. If $q(t)=r(t)e^{i\theta(t)}$ and a
root is chosen initially, its lift is
$z(t)=\sqrt{r(t)}e^{i\theta(t)/2}$. One circuit with an increment
$2\pi$ in $\theta$ ends at $-z(0)$, and two circuits
restore the initial point. This is the transport of the two
lifts over the punctured plane.

An orbital band requires another object: a real line
associated with the orbit and its sign character. Let
$\gamma:S^1\to\mathbb R^2$ be a regular injective
parametrisation of an elliptic orbit, and let $\delta>0$. With
holonomy $-1$ over one circuit, its interval band is
represented by
\begin{equation}
 M_\gamma=([0,2\pi]\times[-\delta,\delta])/
          \bigl((0,u)\sim(2\pi,-u)\bigr).
 \label{vi:vi:eq:mobius-orbital}
\end{equation}
Identifying the base with the image of $\gamma$
places the zero section over the orbit.

\begin{proposition}[Oriented transport over the orbital cycle]
\label{vi:vi:prop:holonomia-orbital}
The line transport in \eqref{vi:vi:eq:mobius-orbital}
satisfies
$\mathcal H(\gamma)=-I$ and $\mathcal H(\gamma^2)=I$.
\end{proposition}
\begin{proof}
The terminal fibre is identified with the initial fibre by
$u\mapsto-u$; composing this map twice gives the identity.
This is the construction in
Theorem~\ref{vi:vi:thm:retorno-mobius}, applied to the orbital loop.
\end{proof}

The covering \eqref{vi:vi:eq:levi-civita} has two-point
fibres; \eqref{vi:vi:eq:mobius-orbital} has interval
fibres. Monodromy of order two appears in both,
while preserving these distinct types. Levi--Civita dynamical
regularisation also uses the transformation
of momenta and time reparametrisation. The present
construction uses only its covering and its lifting
law, with the scope just proved.

\subsection{Phase condition and half-integer shift}
\label{vi:vi:subsec:sommerfeld}

Let $\gamma$ be a cycle of the Hamiltonian realisation, and let
\begin{equation}
 J_\gamma=\oint_\gamma\mathbf p\cdot d\mathbf q
 \label{vi:vi:eq:accion-ciclo}
\end{equation}
be its action. The reduced Planck constant
$\hbar_{\mathrm{int}}>0$ is received from the HMT
action construction, in the same dimensional line as $J_\gamma$.
The quotient $J_\gamma/\hbar_{\mathrm{int}}$ is therefore
dimensionless. Suppose that semiclassical transport in the
complexified line contributes the phase
$\exp(iJ_\gamma/\hbar_{\mathrm{int}})$ and that the real
orientation line contributes the character
$\varepsilon_\gamma\in\{+1,-1\}$. The closure condition
for the transported section is
\begin{equation}
 \varepsilon_\gamma
       \exp\!\left(\frac{iJ_\gamma}{\hbar_{\mathrm{int}}}\right)=1.
 \label{vi:vi:eq:cierre-sommerfeld}
\end{equation}
The product preserves both contributions: the action of the
cycle and the sign of the line. The latter is not absorbed into a
modification of the Planck constant.

\begin{theorem}[Integer and half-integer action sectors]
\label{vi:vi:thm:sommerfeld}
Under the transport \eqref{vi:vi:eq:cierre-sommerfeld},
\begin{equation}
 J_\gamma=
 \begin{cases}
  2\pi\hbar_{\mathrm{int}}n,
          &\varepsilon_\gamma=+1,\\[1mm]
  2\pi\hbar_{\mathrm{int}}(n+\tfrac12),
          &\varepsilon_\gamma=-1,
 \end{cases}
 \qquad n\in\mathbb Z.
 \label{vi:vi:eq:accion-sommerfeld}
\end{equation}
\end{theorem}
\begin{proof}
For a positive sign, the exponential must equal $1$, whose
real arguments are $2\pi n$. For a negative sign, it must
equal $-1$, whose arguments are $(2n+1)\pi$. Multiplying
by $\hbar_{\mathrm{int}}$ produces the two families.
\end{proof}

The formula determines the closure constraint on the
action of a given cycle. Obtaining an atomic spectrum
also uses the specific Hamiltonian, the independent
cycles and the quantisation prescription. If this
prescription includes Maslov indices, their phases are
composed with \eqref{vi:vi:eq:cierre-sommerfeld} according to the
chosen transport, avoiding counting the same
topological contribution twice.

Reversing the orientation of the cycle changes $J_\gamma$ to
$-J_\gamma$, whereas the sign character satisfies
$\varepsilon_{\gamma^{-1}}=\varepsilon_\gamma^{-1}
=\varepsilon_\gamma$. The total phase becomes its
inverse, and the closure condition is preserved. This is a
precise law of orbital bidirectionality. Its identification
with material antisections requires the lift
in Section~\ref{vi:vi:subsec:antisecciones}; equality
of two signs does not by itself establish that identification.

The realization thus brings together three successive
conservations: flux passes from the interior to the boundary;
central dynamics conserves angular momentum and the area
swept per unit time; transport of the section
preserves its closure condition through action and
holonomy. Each result retains the equation, the domain
and the data that support it. The material history of the
TPK is represented in those dynamics with its memory, instead
of being replaced by the orbital figure it publishes.
